\section{CAPÍTULO 2: PRUEBAS, OPTIMIZACIÓN Y DESPLIEGUE}

\subsection{2.1. PRUEBAS UNITARIAS, DE INTEGRACIÓN Y DE ACEPTACIÓN}




\subsubsection{Estrategia general de pruebas}




Para dar estricto cumplimiento a las directrices del Módulo Formativo IV (Implementación de Pruebas de Software), el aseguramiento de la calidad de ParkSys se ejecutó bajo la metodología de la Pirámide de Pruebas de Software, segmentando la verificación en tres niveles formales. La estrategia adoptada consistió en realizar una auditoría endpoint por endpoint tanto en el Backend (API REST con FastAPI) como en el Frontend (Angular 18), registrando para cada funcionalidad la tecnología utilizada, los errores detectados en la primera versión, las correcciones aplicadas en la segunda versión y las recomendaciones pendientes para futuras iteraciones.

\begin{figure}[H]
\caption{Pirámide de pruebas de software aplicada a ParkSys}
\centering
\includegraphics[width=\textwidth,height=9cm,keepaspectratio]{IMAGENES/figuras/figura_16.png}
\vspace{0.2cm}
{\raggedright \textit{Nota.} Representa los tres niveles de la pirámide de pruebas —pruebas unitarias en la base, pruebas de integración en el nivel intermedio y pruebas de aceptación (UAT) en la cúspide— aplicados sobre los módulos billing.py, security.py, schemas.py, la API REST, SQLAlchemy, PostgreSQL y WebSockets. \par}
\end{figure}

\subsubsection{Herramientas y tecnologías utilizadas}




\begin{table}[hbt!]
\caption{Herramientas y tecnologías empleadas en las pruebas de ParkSys}
\label{tab:herramientas_pruebas}
\begin{tabular}{@{}p{3cm} p{1.5cm} p{8cm} p{2.5cm}@{}}
\toprule
\textbf{Herramienta} & \textbf{Versión} & \textbf{Propósito} & \textbf{Capa} \\
\midrule
pytest & 8.x & Framework principal de pruebas en el backend Python & Backend \\
pytest-asyncio & 0.23+ & Soporte para pruebas asíncronas con async/await & Backend \\
httpx & 0.27 & Cliente HTTP asíncrono para pruebas de endpoints FastAPI & Backend \\
aiosqlite & 0.20+ & Driver SQLite asíncrono para base de datos en memoria & Backend \\
SQLAlchemy AsyncSession & 2.0.35 & ORM asíncrono para crear fixtures y verificar persistencia & Backend \\
Jasmine & 5.x & Framework de pruebas para componentes Angular 18 & Frontend \\
Karma & 6.x & Test runner para ejecutar pruebas Jasmine en navegador & Frontend \\
Angular Signals & 18.x & Detección de cambios reactiva para validar re-renders & Frontend \\
RxJS & 7.x & Pruebas de flujos reactivos (WebSocket, Observables) & Frontend \\
\bottomrule
\end{tabular}
\\ \vspace{0.2cm}
\textit{Nota.} Enumera las herramientas de prueba utilizadas en cada capa del sistema, con su versión y propósito técnico.
\end{table}

\subsubsection{Cuadros consolidados de evaluación endpoint por endpoint}




\begin{table}[hbt!]
\caption{Pruebas unitarias y de integración en endpoints — Backend (FastAPI / Python)}
\label{tab:pruebas_backend}
\begin{tabular}{@{}p{2.5cm} p{1.8cm} p{2.2cm} p{3.0cm} p{3.5cm} p{3cm}@{}}
\toprule
\textbf{Endpoint} & \textbf{Tipo prueba} & \textbf{Tech} & \textbf{V1: tiempo/errores} & \textbf{V2: correcciones} & \textbf{Recomendaciones} \\
\midrule
POST /auth/login & Unit / Integr. & pytest, httpx & 120 ms. Error HTTP 500 si usuario no existía. & Manejo HTTP 401. Tiempo: 85 ms. & Rate limiting (5/min). \\
POST /auth/register & Unitaria & pytest, Pydantic & 95 ms. Aceptaba contraseñas cortas y correos inválidos. & Validación estricta con Pydantic. Tiempo: 80 ms. & Verificación correo por token. \\
POST /reservations & Integración & pytest, SQLAlch. & 200 ms. Condición de carrera en reservas concurrentes. & Validación dentro de transacción atómica. Tiempo: 110 ms. & SELECT FOR UPDATE. \\
GET /spaces & Unitaria & pytest, httpx & 150 ms. Consultas N+1 por reserva. & joinedload() e índice B-Tree. Tiempo: 42 ms. & Caché Redis (TTL 5s). \\
POST /payments & Integración & pytest, Stripe & 300 ms. Timeout y errores en montos menores a S/2.00. & Reintentos y validación de monto mínimo. Tiempo: 250 ms. & Webhooks Stripe en 2do plano. \\
POST /vehicles & Unitaria & pytest, asyncio & 140 ms. Aceptaba placas nulas/vacías. & Validación regex estricta (Pydantic). Tiempo: 90 ms. & OCR en garita. \\
PATCH /vehicles/../checkout & Integración & pytest, SQLAlch. & 190 ms. Cálculo incorrecto de duración a medianoche. & math.ceil y timezone.utc consistentes. Tiempo: 18 ms. & Alerta por >24 horas de estadía. \\
GET /reports/pdf/id & Integración & pytest, ReportLab & 350 ms. Desfase UTC de 5 horas. & zoneinfo.ZoneInfo("America/Lima"). Tiempo: 280 ms. & Worker Celery+RabbitMQ para PDF. \\
\bottomrule
\end{tabular}
\\ \vspace{0.2cm}
\textit{Nota.} Consolida la evolución de las pruebas backend en el formato Versión 1 $\rightarrow$ Versión 2 $\rightarrow$ Recomendaciones.
\end{table}

\begin{table}[hbt!]
\caption{Pruebas de consumo de endpoints — Frontend (Angular 18 / TypeScript)}
\label{tab:pruebas_frontend}
\begin{tabular}{@{}p{3.5cm} p{2.0cm} p{2cm} p{2.5cm} p{3cm} p{2.5cm}@{}}
\toprule
\textbf{Componente} & \textbf{Tipo prueba} & \textbf{Tech} & \textbf{V1: tiempo/errores} & \textbf{V2: correcciones} & \textbf{Recomendaciones} \\
\midrule
AuthService $\rightarrow$ login & Unitaria & Jasmine, Karma & 180 ms. UI bloqueada sin mensaje ante credencial errónea. & JwtInterceptor y MatSnackBar. Redirección rol. & Pruebas E2E (Cypress). \\
RegisterComponent & Unitaria & Jasmine, Karma & 200 ms. Sin error en tiempo real. & Validaciones FormGroup reactivas. & Debounce validación correo. \\
NewReservationComponent & Unitaria & Jasmine, Karma & 250 ms. Mapa no refrescaba tras reserva. & Angular Signals y evento WebSocket. & Animaciones (fade-in ámbar). \\
MapComponent & Integración & Jasmine, RxJS & 300 ms. Polling HTTP de alto consumo. & WebSocketService; fin de polling. & Canvas/WebGL para $>1000$ espacios. \\
GaritaPanelComponent & Integración & Jasmine, Karma & 200 ms. Desconexión silenciosa WS. & Reconexión automática (3s) y alerta. & Modo offline (PWA). \\
PaymentComponent & Integración & Jasmine, HttpClient & 280 ms. Doble clic sin spinner. & Deshabilitación y spinner; validación monto. & Integrar Yape/Plin. \\
CheckoutComponent & Integración & Jasmine, Karma & 220 ms. PDF no descargaba auto. & blob y window.open(). & Impresora POS térmica. \\
\bottomrule
\end{tabular}
\\ \vspace{0.2cm}
\textit{Nota.} Consolida la evolución de las pruebas frontend en el formato Versión 1 $\rightarrow$ Versión 2 $\rightarrow$ Recomendaciones.
\end{table}

\begin{table}[hbt!]
\caption{Reporte de errores generales y optimización de rendimiento por módulo}
\label{tab:errores_generales}
\begin{tabular}{@{}p{2.5cm} p{2cm} p{2cm} p{3.5cm} p{3cm} p{3cm}@{}}
\toprule
\textbf{Módulo/Capa} & \textbf{Tipo análisis} & \textbf{Tech} & \textbf{V1: problemas} & \textbf{V2: mejoras} & \textbf{Recomendaciones} \\
\midrule
Base de datos & Rendimiento & PostgreSQL & Consultas lentas por placa, queries O(N). & Índices B-Tree en 10 campos. O(log N). & Particionar vehicle\_logs. \\
Facturación & Error & ReportLab & Desfase UTC en PDF. & zoneinfo America/Lima en backend/front. & PDFs asíncronos. \\
Frontend SPA & Rendimiento & Angular 18 & FCP de 3.8s en móviles. & Standalone Components y Lazy Loading. FCP 1.1s. & Habilitar SSR (Angular Univ.). \\
API REST & Error & FastAPI & CORS bloqueado en Railway. & Regex allow\_origin\_regex. & CSP y restricción CORS. \\
WebSockets & Error & FastAPI+Ang. & Desconexión silenciosa. & Reconexión automática. & Ping/pong heartbeat. \\
Pasarela & Error & Stripe API & HTTP 400 por $<$ S/ 2.00. & Derivación a efectivo. & MercadoPago alternativo. \\
\bottomrule
\end{tabular}
\\ \vspace{0.2cm}
\textit{Nota.} Sintetiza los hallazgos y mejoras de rendimiento y estabilidad a nivel de sistema completo.
\end{table}

\subsubsection{2.1.1. Pruebas unitarias (Unit Testing)}




Las pruebas unitarias verificaron de forma aislada las funciones puras y algoritmos críticos del sistema, sin dependencia de bases de datos externas ni servidores HTTP. Se ejecutaron 30 pruebas unitarias distribuidas en tres componentes fundamentales del backend.

\textbf{Archivo fuente: backend/tests/test\_unit\_billing.py} \\
\textbf{a) Motor de facturación (billing.py) — 14 pruebas}\\
Se verificó que el algoritmo calculase exactamente los minutos efectivos a partir de la diferencia entre fecha\_ingreso y fecha\_salida, redondeando hacia arriba (math.ceil).

\begin{verbatim}
# UT-04 — Redondeo hacia arriba (ceil)
def test_ut04_redondeo_hacia_arriba(self):
    ingreso = datetime(2026, 8, 15, 10, 0, 0, tzinfo=timezone.utc)
    salida = ingreso + timedelta(seconds=90)
    result = calculate_duration_minutes(ingreso, salida)
    assert result == 2
\end{verbatim}

\textbf{b) Seguridad criptográfica (security.py) — 10 pruebas}\\
Se verificó la integridad del hash Argon2id y validación JWT.

\begin{verbatim}
# UT-15 — Hash Argon2 nunca es texto plano
def test_ut15_hash_no_es_texto_plano(self):
    plain = "admin123"
    hashed = hash_password(plain)
    assert hashed != plain
    assert "$argon2" in hashed
\end{verbatim}

\textbf{c) Validación de esquemas Pydantic (schemas.py) — 6 pruebas}
\begin{verbatim}
# UT-26 — Password menor a 8 caracteres
def test_ut26_registro_password_corto(self):
    with pytest.raises(ValidationError):
        RegisterRequest(nombre="Test", correo="t@t.com", password="123")
\end{verbatim}

\begin{table}[hbt!]
\caption{Matriz completa de pruebas unitarias (30 pruebas)}
\label{tab:matriz_ut}
\begin{tabular}{@{}p{2.5cm} p{3.5cm} p{4cm} p{3.5cm} p{2cm}@{}}
\toprule
\textbf{ID} & \textbf{Componente} & \textbf{Entrada} & \textbf{Salida esperada} & \textbf{Resultado} \\
\midrule
UT-01 a UT-08 & billing.py / calculate\_dur & Casos de duración exactos/fracc. & Minutos con redondeo ceil & Conforme \\
UT-09 a UT-14 & billing.py / calculate\_chr & Tipo veh., minutos, tarifa & Monto total calculado & Conforme \\
UT-15 a UT-19 & security.py / hash, verify & Contraseñas prueba & Hash Argon2, validación ok & Conforme \\
UT-20 a UT-24 & security.py / JWT & Payloads acceso/refresco & Tokens gen/decod/rech & Conforme \\
UT-25 a UT-30 & schemas.py & Datos válidos/inválidos & ValidationError o ok & Conforme \\
\bottomrule
\end{tabular}
\\ \vspace{0.2cm}
\textit{Nota.} Resume la matriz completa de 30 pruebas unitarias.
\end{table}

\subsubsection{2.1.2. Pruebas de integración (Integration Testing)}




Se ejecutaron 12 pruebas de integración.

\begin{verbatim}
# IT-02 — Crear reserva cambia espacio a 'reservado'
async def test_it02_crear_reserva_cambia_espacio_a_reservado(...):
    assert sample_space.estado == SpaceStatus.disponible
    # Insert logic...
    assert space_db.estado == SpaceStatus.reservado
\end{verbatim}

\begin{table}[hbt!]
\caption{Matriz completa de pruebas de integración (12 pruebas)}
\label{tab:matriz_it}
\begin{tabular}{@{}p{2cm} p{4cm} p{4cm} p{3cm} p{2cm}@{}}
\toprule
\textbf{ID} & \textbf{Escenario} & \textbf{Procedimiento} & \textbf{Criterio de éxito} & \textbf{Resultado} \\
\midrule
IT-01a/b/c & Autenticación endpoints & Petición sin token/inválido & HTTP 401/403/200 & Conforme \\
IT-02 & Reserva $\rightarrow$ reservado & Crear reserva cajón disp. & Espacio actualizado & Conforme \\
IT-03 & Doble reserva placa & Segunda reserva placa activa & Validación impide dup & Conforme \\
IT-04, 11-12 & Búsquedas indexadas & Consulta código, placa & Búsqueda eficiente & Conforme \\
IT-05 & Cancelación libera espacio & Cancelar reserva pendiente & Espacio disponible & Conforme \\
IT-06, 07 & Transición garita & Walk-in y checkout & Estado cambia ocup/disp & Conforme \\
IT-08 & Cajón bloqueado & reservable=False & HTTP 409 Rechazo & Conforme \\
IT-09, 10 & Bitácora vehicle\_logs & Ingreso/salida walk-in & Log persistido con mins & Conforme \\
\bottomrule
\end{tabular}
\\ \vspace{0.2cm}
\textit{Nota.} Resume la matriz completa de 12 pruebas de integración.
\end{table}

\subsubsection{2.1.3. Pruebas de aceptación (UAT)}

Diseñadas bajo el formato formal Gherkin (Dado / Cuando / Entonces), simulando los flujos completos que los usuarios ejecutarían en producción en el local ParkSys Central Ayacucho. Se ejecutaron 6 pruebas de aceptación.

\textbf{Archivo fuente:} \texttt{backend/tests/test\_acceptance\_flows.py}

\textbf{Ejemplo completo — UAT-03: Cobro exacto al minuto}
\begin{verbatim}
class TestUAT03CobroExacto:
 """
 DADO: Un vehículo permaneció exactamente 1h 15min en el estacionamiento.
 CUANDO: El trabajador presiona "Dar Salida" en el panel.
 ENTONCES: Se liquida hora base + fracción proporcional exacta = S/ 6.25.
 """
 async def test_uat03_cobro_exacto_con_fraccion(self, db_session):
     fecha_ingreso = datetime(2026, 8, 20, 10, 0, 0, tzinfo=timezone.utc)
     fecha_salida = datetime(2026, 8, 20, 11, 15, 0, tzinfo=timezone.utc)
     charge = await calculate_charge(db_session, VehicleType.auto,
                                     fecha_ingreso, fecha_salida)
     assert charge["total"] == 6.25
     assert charge["full_hours"] == 1
     assert charge["remaining_minutes"] == 15
\end{verbatim}





Se ejecutaron 6 pruebas de aceptación (Gherkin).

\begin{table}[hbt!]
\caption{Matriz de pruebas de aceptación (6 pruebas)}
\label{tab:matriz_uat}
\begin{tabular}{@{}p{4cm} p{3.5cm} p{3.5cm} p{4.5cm}@{}}
\toprule
\textbf{Escenario} & \textbf{Condición previa (Dado)} & \textbf{Acción (Cuando)} & \textbf{Resultado (Entonces)} \\
\midrule
UAT-01: Reserva exitosa & Conductor autenticado & Fija fecha/hora y reserva & Código 8 dig, notificación \\
UAT-02: Admisión walk-in & Operador en garita & Digita placa, confirma & Autoasignación, espacio ocupado \\
UAT-03: Cobro exacto al min & Vehículo permaneció 75min & Pulsa "Dar Salida" & 75 min * tarifa = S/6.25, recibo \\
UAT-04: WebSocket push & Operador con WS activo & Conductor crea reserva web & Evento refresco en $<$ 200ms \\
UAT-05: Cancelación reserva & Conductor reserva B-03 & Pulsa "Cancelar" & Cajón disp, reserva cancelada \\
UAT-06: Expiración auto & Reserva vencida & APScheduler ejecuta & Estado vencida, cajón liberado \\
\bottomrule
\end{tabular}
\\ \vspace{0.2cm}
\textit{Nota.} Sintetiza los seis flujos de aceptación validados en formato Gherkin.
\end{table}

\begin{figure}[H]
\caption{Resumen del total de pruebas ejecutadas por nivel}
\centering
\includegraphics[width=\textwidth,height=9cm,keepaspectratio]{IMAGENES/figuras/figura_17.png}
\vspace{0.2cm}
{\raggedright \textit{Nota.} Gráfico de barras que muestra la distribución de las 48 pruebas automatizadas ejecutadas (30 unitarias, 12 de integración y 6 de aceptación), con una tasa de éxito del 100 por ciento en todas ellas. \par}
\end{figure}

\subsection{2.2. REPORTE DE ERRORES}

Durante las etapas de desarrollo, integración y despliegue del sistema se registraron incidencias técnicas que fueron debidamente categorizadas, diagnosticadas y subsanadas. Cada error fue documentado con su identificador, componente afectado, descripción detallada de la falla, nivel de impacto, causa raíz técnica y la solución aplicada.




\begin{table}[hbt!]
\caption{Registro formal de errores identificados y resueltos}
\label{tab:registro_errores_v2}
\begin{tabular}{@{}p{1.2cm} p{2cm} p{4cm} p{1.8cm} p{2.5cm} p{3.5cm}@{}}
\toprule
\textbf{ID} & \textbf{Capa} & \textbf{Falla} & \textbf{Impacto} & \textbf{Causa} & \textbf{Solución} \\
\midrule
ERR-01 & Base de Datos & Condición de carrera: reservas duplicadas simultáneas. & Alto & Sin bloqueo transaccional concurrente. & SELECT ... WHERE estado='disponible' (expire\_on\_commit=False). \\
ERR-02 & WebSockets & Desconexión silenciosa en inactividad. & Medio & Sin reconexión, micro-cortes WiFi. & Reconexión (3s) en Angular, limpieza de sockets backend. \\
ERR-03 & Zonas Horarias & Desfase de 5h UTC en PDF. & Medio & Servidor UTC, sede UTC-5. & zoneinfo.ZoneInfo("America/Lima"). \\
ERR-04 & Stripe API & HTTP 400 por cobro $<$ S/ 2.00. & Medio & Mínimo de Stripe 0.50 USD. & Validación backend: derivar a efectivo. \\
ERR-05 & CORS Deploy & CORS bloqueado en subdominios Railway. & Alto & ALLOWED\_ORIGINS estático. & Regex allow\_origin\_regex en middleware CORS. \\
\bottomrule
\end{tabular}
\\ \vspace{0.2cm}
\textit{Nota.} Registra los cinco errores críticos y de impacto medio.
\end{table}

\begin{figure}[H]
\caption{Flujo de diagnóstico y resolución de errores (ERR-01 a ERR-05)}
\centering
\includegraphics[width=\textwidth,height=9cm,keepaspectratio]{IMAGENES/figuras/figura_18.png}
\vspace{0.2cm}
{\raggedright \textit{Nota.} Diagrama de flujo que ilustra el proceso seguido para cada incidencia. \par}
\end{figure}

\begin{table}[hbt!]
\caption{Distribución de errores por nivel de impacto}
\label{tab:distribucion_errores}
\begin{tabular}{@{}p{6cm} p{3cm} p{6cm}@{}}
\toprule
\textbf{Nivel de impacto} & \textbf{Cantidad} & \textbf{Identificadores} \\
\midrule
Alto (bloqueante para producción) & 2 & ERR-01, ERR-05 \\
Medio (afecta experiencia/precisión) & 3 & ERR-02, ERR-03, ERR-04 \\
Bajo (mejoras de UX sin bloqueo) & 2 & Hallazgos UAT en sección 2.4 \\
\bottomrule
\end{tabular}
\\ \vspace{0.2cm}
\textit{Nota.} Clasifica los errores documentados según su severidad.
\end{table}

\subsection{2.3. OPTIMIZACIÓN DEL RENDIMIENTO}

Para asegurar que ParkSys mantenga tiempos de respuesta mínimos bajo condiciones de alta concurrencia en horas punta, se ejecutaron optimizaciones de ingeniería en dos capas: Frontend y Backend.




\subsubsection{2.3.1. Frontend: Lazy Loading y Standalone Components}

Angular 18 permite definir cada componente como Standalone (standalone: true), eliminando la necesidad de NgModules tradicionales. ParkSys aprovecha esta capacidad junto con Lazy Loading para reducir drásticamente el peso inicial de la aplicación. Los módulos de administrador, trabajador y conductor se cargan solo cuando el usuario navega a su sección correspondiente mediante \texttt{loadComponent()} en el archivo \texttt{app.routes.ts}.




\begin{table}[hbt!]
\caption{Componentes con Lazy Loading implementado}
\label{tab:lazy_loading_v2}
\begin{tabular}{@{}p{3cm} p{9cm} p{3cm}@{}}
\toprule
\textbf{Módulo} & \textbf{Componentes Lazy-Loaded} & \textbf{Ahorro est.} \\
\midrule
Administrador & Financials, Personnel, Settings, Support, Audit, Clients, AdminReviews, TransactionDetail & $\sim$ 180 KB \\
Conductor & SecuritySettings, NotificationSettings, EditProfile, Receipt, PendingPayments, PaymentSuccess & $\sim$ 120 KB \\
Público & ResetPassword, Terms, PublicLegal, PublicSupport, PublicPaymentSuccess & $\sim$ 60 KB \\
\bottomrule
\end{tabular}
\\ \vspace{0.2cm}
\textit{Nota.} Detalla la distribución de componentes cargados de forma diferida.
\end{table}

\begin{figure}[H]
\caption{Comparativa de First Contentful Paint antes y después de la optimización}
\centering
\includegraphics[width=\textwidth,height=9cm,keepaspectratio]{IMAGENES/figuras/figura_19.png}
\vspace{0.2cm}
{\raggedright \textit{Nota.} Gráfico comparativo que muestra la reducción del FCP de 3.8 segundos a 1.1 segundos. \par}
\end{figure}

\subsubsection{2.3.2. Backend y base de datos: asyncpg e indexación}

\textbf{Driver asíncrono asyncpg}

El uso del driver nativo asyncpg con SQLAlchemy 2.0 Async permite que FastAPI atienda múltiples peticiones concurrentes sin bloquear el hilo principal esperando respuestas de disco:

\begin{verbatim}
# backend/app/db/session.py
engine = create_async_engine(
 settings.database_url, # postgresql+asyncpg://...
 echo=settings.environment == "development",
 pool_pre_ping=True,
)

async_session_factory = async_sessionmaker(
 engine,
 class_=AsyncSession,
 expire_on_commit=False,
)
\end{verbatim}

\textbf{Resultado:} el servidor soporta más de 800 peticiones por segundo en benchmarks locales con uvicorn y el pool de conexiones de asyncpg.

\textbf{Índices B-Tree en campos de búsqueda frecuente}

Se crearon 10 índices explícitos en las columnas más consultadas, reduciendo la complejidad de búsqueda de O(N) a O(log N).




\begin{table}[hbt!]
\caption{Campos indexados en el sistema}
\label{tab:indices_bd_v2}
\begin{tabular}{@{}p{3cm} p{3.5cm} p{3.5cm} p{5cm}@{}}
\toprule
\textbf{Tabla} & \textbf{Campo} & \textbf{Tipo de índice} & \textbf{Uso principal} \\
\midrule
users & correo & UNIQUE & Login y registro \\
spaces & codigo & UNIQUE & Búsqueda de cajón \\
reservations & codigo\_reserva & UNIQUE & Búsqueda en garita \\
reservations & usuario\_id & B-Tree FK & Listado por conductor \\
reservations & espacio\_id & B-Tree FK & Verificación de disponibilidad \\
vehicle\_logs & placa & B-Tree & Búsqueda por placa \\
vehicle\_logs & reserva\_id & B-Tree FK & Join reservations \\
vehicle\_logs & espacio\_id & B-Tree FK & Historial por espacio \\
payments & vehicle\_log\_id & B-Tree FK & Asociación pago-registro \\
payments & reserva\_id & B-Tree FK & Asociación pago-reserva \\
\bottomrule
\end{tabular}
\\ \vspace{0.2cm}
\textit{Nota.} Enumera los diez índices creados sobre columnas de mayor frecuencia.
\end{table}

\begin{table}[hbt!]
\caption{Métricas de optimización y benchmarks}
\label{tab:metricas_optimizacion}
\begin{tabular}{@{}p{5cm} p{3cm} p{3cm} p{4cm}@{}}
\toprule
\textbf{Indicador} & \textbf{Valor base} & \textbf{Valor opt.} & \textbf{Mejora} \\
\midrule
First Contentful Paint (FCP) & 3.8 s & 1.1 s & 71\% más rápido \\
Latencia /spaces/availability & 280 ms & 42 ms & 85\% de reducción \\
Tiempo de cálculo de checkout & 190 ms & 18 ms & 90\% de reducción \\
Consumo de RAM en reposo & 180 MB & 78 MB & 56\% más liviano \\
Throughput del servidor & 320 req/s & 800+ req/s & 150\% de incremento \\
\bottomrule
\end{tabular}
\\ \vspace{0.2cm}
\textit{Nota.} Sintetiza las métricas de rendimiento antes y después.
\end{table}

\begin{figure}[H]
\caption{Arquitectura de indexación y acceso asíncrono a la base de datos}
\centering
\includegraphics[width=\textwidth,height=9cm,keepaspectratio]{IMAGENES/figuras/figura_20.png}
\vspace{0.2cm}
{\raggedright \textit{Nota.} Diagrama que ilustra el flujo con asyncpg e índices B-Tree. \par}
\end{figure}

\begin{figure}[H]
\caption{Comparativa de throughput del servidor antes y después de la optimización}
\centering
\includegraphics[width=\textwidth,height=9cm,keepaspectratio]{IMAGENES/figuras/figura_21.png}
\vspace{0.2cm}
{\raggedright \textit{Nota.} Gráfico de barras que compara el throughput del servidor. \par}
\end{figure}

\subsection{2.4. PRUEBAS DE USUARIO (UAT)}

Las pruebas de usuario se ejecutaron en un entorno controlado con usuarios reales representando los tres roles del sistema: Administrador, Trabajador de garita y Conductor. El objetivo fue validar que los flujos de trabajo del sistema se ejecutaran de manera intuitiva, eficiente y libre de errores desde la perspectiva del usuario final, y no exclusivamente desde la del desarrollador.




\begin{table}[hbt!]
\caption{Parámetros de ejecución de las pruebas de usuario}
\label{tab:parametros_uat}
\begin{tabular}{@{}p{4cm} p{11cm}@{}}
\toprule
\textbf{Parámetro} & \textbf{Detalle} \\
\midrule
Fecha de ejecución & Agosto del 2026 \\
Lugar & Sede ParkSys Central — Ayacucho \\
Participantes & Tres usuarios (un administrador, un trabajador y un conductor) \\
Dispositivos empleados & PC de escritorio (garita), laptop (admin), smartphone Android \\
Entorno de prueba & Producción (Railway), conectado a PostgreSQL remoto \\
Metodología aplicada & Observación directa, cuestionario post-uso y registro de incidencias \\
\bottomrule
\end{tabular}
\\ \vspace{0.2cm}
\textit{Nota.} Describe las condiciones de ejecución de las pruebas.
\end{table}

\begin{figure}[H]
\caption{Evidencia fotográfica de la sesión de pruebas de usuario en la sede ParkSys Central}
\centering
\includegraphics[width=\textwidth,height=9cm,keepaspectratio]{IMAGENES/figuras/figura_22.png}
\vspace{0.2cm}
{\raggedright \textit{Nota.} Registra el momento de la ejecución de las pruebas con usuarios operando el sistema. \par}
\end{figure}

\begin{table}[hbt!]
\caption{Escenarios de prueba — Rol Conductor}
\label{tab:uat_conductor}
\begin{tabular}{@{}p{0.5cm} p{4.5cm} p{4.5cm} p{1.8cm} p{3.7cm}@{}}
\toprule
\textbf{N.°} & \textbf{Tarea} & \textbf{Criterio de éxito} & \textbf{Result.} & \textbf{Observaciones} \\
\midrule
1 & Registro correo/pass & Cuenta creada e inicio & Cumplido & 45 segundos \\
2 & Agregar vehículo & Listado en "Mis Vehículos" & Cumplido & Sin dificultad \\
3 & Crear reserva mapa 2D & Código y cajón en ámbar & Cumplido & Muy intuitivo \\
4 & Consultar estado & Estado "pendiente" & Cumplido & - \\
5 & Cancelar reserva & Estado "cancelada" & Cumplido & Sugirió confirmación UI \\
6 & Pagar Stripe online & Checkout exitoso/recibo & Cumplido & Redirect de 3 seg \\
\bottomrule
\end{tabular}
\\ \vspace{0.2cm}
\textit{Nota.} Detalla las tareas ejecutadas por el conductor.
\end{table}

\begin{table}[hbt!]
\caption{Escenarios de prueba — Rol Trabajador de Garita}
\label{tab:uat_trabajador}
\begin{tabular}{@{}p{0.5cm} p{3.7cm} p{4cm} p{1.8cm} p{4.5cm}@{}}
\toprule
\textbf{N.°} & \textbf{Tarea} & \textbf{Criterio de éxito} & \textbf{Result.} & \textbf{Observaciones} \\
\midrule
1 & Iniciar sesión panel & Mapa en tiempo real & Cumplido & Carga en 1.2s \\
2 & Ingreso walk-in & Espacio/hora registrada & Cumplido & $<$ 10 seg \\
3 & Dar salida (checkout) & Monto mostrado, PDF & Cumplido & PDF automático \\
4 & Verificación mapa WS & Espacio cambia color auto & Cumplido & $<$ 200 ms via WebSocket \\
5 & Buscar por placa & Resultado rápido & Cumplido & 40 ms indexado \\
\bottomrule
\end{tabular}
\\ \vspace{0.2cm}
\textit{Nota.} Describe las tareas del trabajador de garita.
\end{table}

\begin{table}[hbt!]
\caption{Escenarios de prueba — Rol Administrador}
\label{tab:uat_admin}
\begin{tabular}{@{}p{0.5cm} p{3.7cm} p{4cm} p{1.8cm} p{4.5cm}@{}}
\toprule
\textbf{N.°} & \textbf{Tarea} & \textbf{Criterio de éxito} & \textbf{Result.} & \textbf{Observaciones} \\
\midrule
1 & Acceder dashboard KPIs & Indicadores visibles & Cumplido & Vista consolidada OK \\
2 & Gestionar espacios & CRUD cajones & Cumplido & - \\
3 & Gestionar personal & CRUD trabajadores & Cumplido & - \\
4 & Configurar tarifas & Tarifas reflejadas & Cumplido & Verificó cobro moto \\
5 & Reportes financieros & Filtrados por fechas & Cumplido & - \\
6 & Revisar auditoría & Registro temporal OK & Cumplido & - \\
\bottomrule
\end{tabular}
\\ \vspace{0.2cm}
\textit{Nota.} Detalla las tareas administrativas del administrador.
\end{table}

\begin{table}[hbt!]
\caption{Resultados consolidados de las pruebas de usuario}
\label{tab:resultados_uat_v2}
\begin{tabular}{@{}p{8cm} p{7cm}@{}}
\toprule
\textbf{Métrica} & \textbf{Resultado} \\
\midrule
Tasa de completitud de tareas & 100\% (17 de 17 tareas completadas) \\
Errores críticos durante pruebas & 0 \\
Errores menores (sugerencias UX) & 2 (Confirmación cancelación, auto-mayúsculas) \\
Satisfacción general (escala 1-5) & 4.7 \\
Tiempo promedio aprendizaje & Menos de 3 minutos por rol \\
\bottomrule
\end{tabular}
\\ \vspace{0.2cm}
\textit{Nota.} Sintetiza los indicadores cuantitativos de UAT.
\end{table}

\subsection{2.5. DOCUMENTACIÓN TÉCNICA Y MANUAL DE USUARIO}

\subsubsection{2.5.1. Documentación Técnica (Dirigida al Desarrollador)}

La documentación técnica de ParkSys se encuentra distribuida en los siguientes artefactos del repositorio del proyecto:




\begin{table}[hbt!]
\caption{Documentos técnicos del repositorio}
\label{tab:documentos_tecnicos}
\begin{tabular}{@{}p{4cm} p{4cm} p{7.5cm}@{}}
\toprule
\textbf{Documento} & \textbf{Ubicación} & \textbf{Contenido} \\
\midrule
README.md & ./README.md & Descripción, stack, instalación \\
CAPITULO\_1.md & ./CAPITULO\_1.md & Análisis, modelos, arquitectura API \\
CAPITULO\_2\_PRUEBAS & ./CAP...\_PRUEBAS.md & Unit/Integration/UAT, errores \\
DESCRIPCION.md & ./DESCRIPCION.md & Requisitos funcionales \\
.env.example & ./backend/.env... & Plantilla de entorno \\
pyproject.toml & ./backend/py... & Config pytest asíncrono \\
docker-compose.yml & ./docker-compose... & Orquestación local (PG + FastAPI) \\
Dockerfile & ./backend/Docker... & Imagen Python Uvicorn \\
Swagger OpenAPI & /docs (FastAPI) & Documentación interactiva de API \\
\bottomrule
\end{tabular}
\\ \vspace{0.2cm}
\textit{Nota.} Enumera artefactos documentales.
\end{table}

\begin{figure}[H]
\caption{Arquitectura del sistema (Resumen técnico de capas)}
\centering
\includegraphics[width=\textwidth,height=9cm,keepaspectratio]{IMAGENES/figuras/figura_23.png}
\vspace{0.2cm}
{\raggedright \textit{Nota.} Ilustra la disposición en capas del sistema: el frontend en Angular 18 con sus tres módulos por rol, el backend en FastAPI con los routers, la capa de servicios y el ORM asíncrono, y la base de datos PostgreSQL 16 con sus índices B-Tree. \par}
\end{figure}

\begin{table}[hbt!]
\caption{Endpoints de la API REST (15 routers registrados)}
\label{tab:endpoints_api}
\begin{tabular}{@{}p{3.5cm} p{3.5cm} p{8cm}@{}}
\toprule
\textbf{Prefijo} & \textbf{Router} & \textbf{Descripción} \\
\midrule
/api/v1/auth & auth.py & Login, registro, token refresh, OAuth \\
/api/v1/spaces & spaces.py & CRUD cajones, 2D mapa \\
/api/v1/reservations & reservations.py & Reservas, cancelación \\
/api/v1/vehicles & vehicles.py & Walk-in, checkout \\
/api/v1/payments & payments.py & Stripe checkout, transacciones \\
/api/v1/reviews & reviews.py & Reseñas conductores \\
/api/v1/notifications & notifications.py & Alertas, mensajes \\
/api/v1/reports & reports.py & Reportes PDF (ReportLab) \\
/api/v1/settings & settings.py & Tarifas, horarios, tolerancias \\
/api/v1/workers & workers.py & Personal garita \\
/api/v1/history & history.py & Historial log vehicular \\
/api/v1/terms & terms.py & Términos legales \\
/api/v1/audit & audit.py & Bitácora inmutable \\
/api/v1/users & users.py & CRUD perfiles \\
/api/v1/ws & websockets.py & Real-time WebSockets \\
\bottomrule
\end{tabular}
\\ \vspace{0.2cm}
\textit{Nota.} Conjunto de controladores REST de la API.
\end{table}

\begin{table}[hbt!]
\caption{Modelos de base de datos referenciados en la documentación}
\label{tab:modelos_bd}
\begin{tabular}{@{}p{3cm} p{3cm} p{9cm}@{}}
\toprule
\textbf{Modelo} & \textbf{Tabla} & \textbf{Campos principales} \\
\midrule
User & users & id, nombre, correo, password\_hash, rol \\
Space & spaces & id, codigo, tipo\_vehiculo, estado, reservable... \\
Reservation & reservations & id, usuario\_id, espacio\_id, placa, estado \\
VehicleLog & vehicle\_logs & id, placa, ingreso, salida \\
Payment & payments & id, vehicle\_log\_id, monto, estado \\
Rate & rates & id, tipo\_vehiculo, precio\_hora \\
ParkingInfo & parking\_info & id, horarios, nombre, direccion \\
\bottomrule
\end{tabular}
\\ \vspace{0.2cm}
\textit{Nota.} Referencia rápida de ORMs.
\end{table}

\subsubsection{2.5.2. Manual de Usuario (Dirigido al Usuario Final)}

El manual de usuario fue diseñado para los tres perfiles del sistema, con instrucciones secuenciales para cada flujo de trabajo principal.

\begin{table}[hbt!]
\caption{Manual del Conductor}
\label{tab:manual_conductor_v2}
\begin{tabular}{@{}p{0.5cm} p{10.5cm} p{4cm}@{}}
\toprule
\textbf{Paso} & \textbf{Acción} & \textbf{Pantalla} \\
\midrule
1 & Ingresar a la URL del sistema web & Landing Page \\
2 & Pulsar "Registrarse" y completar formulario & Registro \\
3 & Iniciar sesión con correo y contraseña & Login \\
4 & Ingresar a "Mis Vehículos" y agregar & Vehículos \\
5 & Ingresar a "Nueva Reserva" y seleccionar cajón 2D & Mapa 2D \\
6 & Fijar fecha/hora y reservar & Reserva \\
7 & Anotar código alfanumérico & Confirmación \\
8 & Consultar estado en "Mis Reservas" & Listado \\
9 & Pago Stripe online / Efectivo & Pago \\
10 & Descargar comprobante PDF & Historial \\
\bottomrule
\end{tabular}
\\ \vspace{0.2cm}
\textit{Nota.} Procedimiento paso a paso para el conductor.
\end{table}

\begin{table}[hbt!]
\caption{Manual del Trabajador de Garita}
\label{tab:manual_trabajador_v2}
\begin{tabular}{@{}p{0.5cm} p{10.5cm} p{4cm}@{}}
\toprule
\textbf{Paso} & \textbf{Acción} & \textbf{Pantalla} \\
\midrule
1 & Iniciar sesión & Login \\
2 & Visualizar panel en tiempo real & Panel Garita \\
3 & Ingreso con reserva (validar código 8 dígitos) & Búsqueda \\
4 & Ingreso walk-in (digitar placa, auto-asignar) & Form walk-in \\
5 & Dar salida por placa (Checkout autom.) & Checkout \\
6 & Registrar pago (efectivo/tarjeta) & Registro pago \\
7 & Comprobante PDF entregado & Comprobante \\
\bottomrule
\end{tabular}
\\ \vspace{0.2cm}
\textit{Nota.} Operativa en garita.
\end{table}

\begin{table}[hbt!]
\caption{Manual del Administrador}
\label{tab:manual_admin_v2}
\begin{tabular}{@{}p{0.5cm} p{10.5cm} p{4cm}@{}}
\toprule
\textbf{Paso} & \textbf{Acción} & \textbf{Pantalla} \\
\midrule
1 & Iniciar sesión admin & Login \\
2 & Visualizar KPIs en tiempo real & Dashboard \\
3 & CRUD de cajones & Gestor espacios \\
4 & CRUD personal trabajador & Personal \\
5 & Tarifas, horarios, tolerancias & Configuración \\
6 & Reportes PDF de caja & Finanzas \\
7 & Auditoría temporal de eventos & Auditoría \\
8 & Listado clientes & Clientes \\
\bottomrule
\end{tabular}
\\ \vspace{0.2cm}
\textit{Nota.} Funciones administrativas.
\end{table}

\subsection{2.6. DESPLIEGUE DEL SISTEMA}




\begin{figure}[H]
\caption{Arquitectura de despliegue de ParkSys (entorno de producción y entorno local)}
\centering
\includegraphics[width=\textwidth,height=9cm,keepaspectratio]{IMAGENES/figuras/figura_24.png}
\vspace{0.2cm}
{\raggedright \textit{Nota.} Ilustra la disposición de los servicios en producción y local. \par}
\end{figure}

\begin{table}[hbt!]
\caption{Configuración del servicio de base de datos (entorno local)}
\label{tab:config_bd_local}
\begin{tabular}{@{}p{4cm} p{11cm}@{}}
\toprule
\textbf{Parámetro} & \textbf{Valor} \\
\midrule
Imagen & postgres:16-alpine \\
Base de datos & parksys \\
Puerto expuesto & 5432 \\
Volumen persistente & postgres\_data \\
Healthcheck & pg\_isready, cada 5 segundos, 5 reintentos \\
\bottomrule
\end{tabular}
\\ \vspace{0.2cm}
\textit{Nota.} Contenedor PostgreSQL local.
\end{table}

\begin{table}[hbt!]
\caption{Configuración del servicio backend (entorno local)}
\label{tab:config_backend_local}
\begin{tabular}{@{}p{4cm} p{11cm}@{}}
\toprule
\textbf{Parámetro} & \textbf{Valor} \\
\midrule
Build & ./backend/Dockerfile \\
Imagen base & python:3.11-slim \\
Puerto & 8000 \\
Dependencia & Espera healthcheck PostgreSQL \\
Comando & alembic upgrade head \&\& uvicorn app.main:app \\
\bottomrule
\end{tabular}
\\ \vspace{0.2cm}
\textit{Nota.} Contenedor FastAPI local.
\end{table}

\begin{table}[hbt!]
\caption{Criterios de selección de la plataforma de despliegue (Railway)}
\label{tab:criterios_railway}
\begin{tabular}{@{}p{3.5cm} p{12cm}@{}}
\toprule
\textbf{Criterio} & \textbf{Justificación} \\
\midrule
Costo & Capa gratuita MVP \\
Facilidad & CI/CD automático con GitHub Webhook \\
PostgreSQL & Soporte nativo manejado con backups \\
Docker & Soporte nativo de Dockerfile \\
Certificados & HTTPS automático let's encrypt \\
Variables entorno & Panel unificado amigable \\
\bottomrule
\end{tabular}
\\ \vspace{0.2cm}
\textit{Nota.} Justifica Railway como PaaS productivo.
\end{table}

\begin{table}[hbt!]
\caption{Configuración del backend en Railway}
\label{tab:config_railway_backend}
\begin{tabular}{@{}p{4cm} p{11cm}@{}}
\toprule
\textbf{Parámetro} & \textbf{Valor} \\
\midrule
Directorio raíz & /backend \\
Comando Build & pip install -r requirements.txt \\
Comando Inicio & alembic upgrade head \&\& uvicorn ... --port \$PORT \\
Uso Dockerfile & Sí \\
\bottomrule
\end{tabular}
\\ \vspace{0.2cm}
\textit{Nota.} Backend productivo.
\end{table}

\begin{table}[hbt!]
\caption{Configuración del frontend en Railway}
\label{tab:config_railway_frontend}
\begin{tabular}{@{}p{4cm} p{11cm}@{}}
\toprule
\textbf{Parámetro} & \textbf{Valor} \\
\midrule
Directorio raíz & /frontend \\
Comando Build & npm install \&\& ng build --configuration production \\
Dir de salida & dist/frontend/browser \\
Tipo de servicio & Sitio estático SPA \\
\bottomrule
\end{tabular}
\\ \vspace{0.2cm}
\textit{Nota.} Frontend productivo SPA.
\end{table}

\begin{figure}[H]
\caption{Flujo del pipeline de despliegue continuo (CI/CD) de ParkSys}
\centering
\includegraphics[width=\textwidth,height=9cm,keepaspectratio]{IMAGENES/figuras/figura_25.png}
\vspace{0.2cm}
{\raggedright \textit{Nota.} Ilustra la secuencia desde GitHub hasta producción automática. \par}
\end{figure}

\begin{table}[hbt!]
\caption{Variables de entorno configuradas en el entorno de producción}
\label{tab:env_produccion}
\begin{tabular}{@{}p{6cm} p{9cm}@{}}
\toprule
\textbf{Variable} & \textbf{Descripción} \\
\midrule
DATABASE\_URL & Conexión PG (Railway inject) \\
JWT\_SECRET & Firma HS256 JWT \\
STRIPE\_SECRET\_KEY & Pasarela de pagos secreta \\
ALLOWED\_ORIGINS & Dominios CORS permitidos \\
SMTP\_USER & Credenciales correo alertas \\
\bottomrule
\end{tabular}
\\ \vspace{0.2cm}
\textit{Nota.} Sintetiza el .env productivo.
\end{table}

\subsection{2.7. ENTREGA Y PUESTA EN MARCHA}




\begin{table}[hbt!]
\caption{Checklist de entregables del proyecto ParkSys}
\label{tab:checklist_entregables}
\begin{tabular}{@{}p{1cm} p{4cm} p{7cm} p{2cm}@{}}
\toprule
\textbf{N.°} & \textbf{Entregable} & \textbf{Descripción} & \textbf{Estado} \\
\midrule
1 & Código fuente & Repositorio Git completo & Entregado \\
2 & Base de datos & PG 16 migrado + seed & Entregado \\
3 & Backend desplegado & FastAPI activo en Railway & Operativo \\
4 & Frontend desplegado & Angular 18 activo & Operativo \\
5 & Documentación & Manuales técnicos MD & Entregado \\
6 & Manual Usuario & Manuales por perfil & Entregado \\
7 & Seed Data & seed.py ejecutado & Ejecutado \\
8 & Credenciales & Cuentas admin/trab/cond & Entregado \\
9 & Suite de pruebas & 48 automatizadas (100\% ok) & Entregado \\
10 & Docker local & Config docker-compose & Entregado \\
\bottomrule
\end{tabular}
\\ \vspace{0.2cm}
\textit{Nota.} Consolida el estado de cumplimiento entregables.
\end{table}

\begin{table}[hbt!]
\caption{Fases del procedimiento de puesta en marcha}
\label{tab:fases_puesta_marcha}
\begin{tabular}{@{}p{0.5cm} p{9.5cm} p{2cm} p{2cm}@{}}
\toprule
\textbf{Fase} & \textbf{Actividad} & \textbf{Respon.} & \textbf{Estado} \\
\midrule
1 & Verificación infraestructura red garita & Técnico & Cumplido \\
2 & Despliegue producción Railway & Dev & Cumplido \\
3 & Migración DDL alembic & Dev & Cumplido \\
4 & Carga inicial seed.py & Dev & Cumplido \\
5 & Configuración tarifas/tolerancia admin & Admin & Cumplido \\
6 & Diseño mapa 2D interactivo & Admin & Cumplido \\
7 & Capacitación de garita (walk-in/checkout) & Dev/Trabaj & Cumplido \\
8 & Prueba completa real E2E & Todos & Cumplido \\
9 & Acta Aceptación & Admin & Cumplido \\
\bottomrule
\end{tabular}
\\ \vspace{0.2cm}
\textit{Nota.} Describe las 9 fases de puesta en marcha operativa.
\end{table}

\begin{figure}[H]
\caption{Diagrama de flujo del procedimiento de puesta en marcha del sistema ParkSys}
\centering
\includegraphics[width=\textwidth,height=9cm,keepaspectratio]{IMAGENES/figuras/figura_26.png}
\vspace{0.2cm}
{\raggedright \textit{Nota.} Ilustra la secuencia cronológica de puesta en marcha. \par}
\end{figure}

\begin{table}[hbt!]
\caption{Plan de contingencia post-entrega}
\label{tab:plan_contingencia}
\begin{tabular}{@{}p{4cm} p{8cm} p{3cm}@{}}
\toprule
\textbf{Riesgo} & \textbf{Acción} & \textbf{Responsable} \\
\midrule
Caída servidor Railway & Auto-restart; o levantar Docker local. & Dev \\
Pérdida internet garita & Log manual temporal (libreta), sync posterior. & Trabajador \\
Error Base de datos & Restaurar backup diario Railway automático. & Dev \\
Bug productivo & Hotfix git push $\rightarrow$ auto-deploy. & Dev \\
Aumento de tráfico & Scale up de memoria/CPU en Railway panel. & Dev \\
\bottomrule
\end{tabular}
\\ \vspace{0.2cm}
\textit{Nota.} Acciones previstas ante escenarios operativos.
\end{table}

\begin{table}[hbt!]
\caption{Acta de aceptación final del proyecto ParkSys}
\label{tab:acta_aceptacion}
\begin{tabular}{@{}p{4cm} p{11cm}@{}}
\toprule
\textbf{Campo} & \textbf{Detalle} \\
\midrule
Proyecto & ParkSys — Sistema de Gestión Estacionamiento \\
Cliente & Estacionamiento Central — Ayacucho \\
Fecha entrega & Agosto 2026 \\
Equipo & Aspur Yaranga, Condori Cotaquispe \\
Módulos entregados & Conductor, Garita, Admin, Stripe, WS, PDF, Aud \\
Pruebas ejecutadas & 48 automatizadas, 17 tareas reales \\
Tasa de éxito & 100\% ok \\
Conformidad & Aceptado sin observaciones \\
\bottomrule
\end{tabular}
\\ \vspace{0.2cm}
\textit{Nota.} Documenta formalmente la conformidad del cliente.
\end{table}
